\section*{第 \ref{ch:g12:stretch-reflex} 章——牵张反射与神经信息}

\begin{solution}{exo:g12:stretch-reflex:1}
感受器（肌梭）、感觉\omterm{def:g12:stretch-reflex:neuron}{神经元}、整合中心（脊髓灰质，
一次\omterm{def:g12:stretch-reflex:synapse}{突触}）、运动\omterm{def:g12:stretch-reflex:neuron}{神经元}、效应器（肌肉）。
\end{solution}

\begin{solution}{exo:g12:stretch-reflex:2}
\omterm{def:g10:cells-common-unit:cell}{细胞}体含\omterm{def:g10:cells-common-unit:organelle}{细胞核}，树突接收信号，并有一条
轴突（常有髓鞘），负责将信息送至末梢。
信息由树突和\omterm{def:g10:cells-common-unit:cell}{细胞}体沿轴突
传至末梢。
\end{solution}

\begin{solution}{exo:g12:stretch-reflex:3}
膜电位短暂逆转，由 \qty{-70}{mV} 变为约
\qty{+30}{mV}，并在一毫秒内恢复，沿轴突传导时幅度不变。
具有“全或无”特性：低于阈值时无反应，高于阈值时总是产生
相同的完整信号。强度由\omterm{prop:g12:stretch-reflex:actionpotential}{动作电位}的
频率编码。
\end{solution}

\begin{solution}{exo:g12:stretch-reflex:4}
\omterm{prop:g12:stretch-reflex:actionpotential}{动作电位}到达末梢；\omterm{def:g12:stretch-reflex:synapse}{突触}小泡将\omterm{def:g12:stretch-reflex:synapse}{神经递质}释放
至\omterm{def:g12:stretch-reflex:synapse}{突触}间隙；递质与接收\omterm{def:g10:cells-common-unit:cell}{细胞}上的受体结合，
使其膜电位发生改变（趋近或远离阈值）；
随后递质被分解或重新摄取。
\end{solution}

\begin{solution}{exo:g12:stretch-reflex:5}
乙酰胆碱。箭毒占据肌肉上的受体，但不激活它们；
肉毒毒素则阻止乙酰胆碱从小泡中释放。
两者均导致麻痹。
\end{solution}

\begin{solution}{exo:g12:stretch-reflex:6}
\qty{100}{ms} 内有 6、12 和 24 个峰电位，即每秒 60、120 和 240 个。
峰电位的大小和形状保持不变。
\end{solution}

\begin{solution}{exo:g12:stretch-reflex:7}
传导时间 $30 - 12 = \qty{18}{ms}$，对应 \qty{1.5}{m}，约为
\qty{83}{m/s}。
\end{solution}

\begin{solution}{exo:g12:stretch-reflex:8}
背根携带感觉纤维：剪断后，牵张信息无法传入脊髓，
但运动\omterm{def:g12:stretch-reflex:neuron}{神经元}及其纤维完好，仍可受刺激而使肌肉收缩。
腹根携带运动纤维：剪断后，
脊髓的信号完全无法到达肌肉，
肌肉得不到任何来自脊髓的指令。
\end{solution}

\begin{solution}{exo:g12:stretch-reflex:9}
若拮抗肌也收缩，它将阻碍运动并
被拉伸，从而引发它自己的\omterm{def:g12:stretch-reflex:reflex}{反射}——一场拔河。
脊髓中的中间\omterm{def:g12:stretch-reflex:neuron}{神经元}起抑制作用，经抑制性
\omterm{def:g12:stretch-reflex:synapse}{突触}作用于拮抗肌的运动\omterm{def:g12:stretch-reflex:neuron}{神经元}。
\end{solution}

\begin{solution}{exo:g12:stretch-reflex:10}
1 个单位：无反应（低于阈值）。4 和 8 个单位：均产生\omterm{prop:g12:stretch-reflex:actionpotential}{动作电位}，
幅度相同，均为 \qty{100}{mV}，但频率较高，
且 8 个单位时的频率高于 4 个单位时。
\end{solution}

\begin{solution}{exo:g12:stretch-reflex:11}
只有末梢含有递质小泡，只有接收膜
含有受体，因此化学信号只能单向传递。递质的释放、
递质在\omterm{def:g12:stretch-reflex:synapse}{突触}间隙中的扩散以及受体的反应，
大约需要一毫秒。
\end{solution}

\begin{solution}{exo:g12:stretch-reflex:12}
乙酰胆碱在连接处积累，使纤维持续放电：
先是抽搐，继而肌肉持续收缩直至耗竭，
包括呼吸肌——最终窒息死亡。
受体阻断剂（一种类阿托品药物）可减轻过度
刺激；剂量过少则中毒不解，过多则如箭毒般引起麻痹，
因此剂量必须恰当。
\end{solution}

\begin{solution}{exo:g12:stretch-reflex:13}
每一种感觉输入都兴奋其运动\omterm{def:g12:stretch-reflex:neuron}{神经元}，通常还经中间\omterm{def:g12:stretch-reflex:neuron}{神经元}抑制
拮抗肌的运动\omterm{def:g12:stretch-reflex:neuron}{神经元}。若抑制被阻断，一次触觉刺激
会同时兴奋每个\omterm{def:g10:muscles-and-joints:joint}{关节}的两侧，兴奋不受制约地
在脊髓中扩散：所有肌肉同时收缩，
痉挛又拉伸其他肌肉，它们的\omterm{def:g12:stretch-reflex:reflex}{反射}随之叠加。
\end{solution}

\begin{solution}{exo:g12:stretch-reflex:14}
传导时间 $60 - 12 = \qty{48}{ms}$，对应 \qty{1.5}{m}，约为
\qty{31}{m/s}，仅为正常速度的三分之一。髓鞘使\omterm{prop:g12:stretch-reflex:actionpotential}{动作电位}
在郎飞结之间跳跃式传导，而非连续传导；
若无髓鞘，纤维的传导方式便如同细的无髓纤维。
\end{solution}

\begin{solution}{exo:g12:stretch-reflex:15}
脊髓在脑与身体之间传递信息，同时还对信息进行处理：
其灰质含有\omterm{def:g12:stretch-reflex:reflex}{反射弧}的\omterm{def:g12:stretch-reflex:synapse}{突触}，在此，牵张被转化为收缩，
并经中间\omterm{def:g12:stretch-reflex:neuron}{神经元}转化为对拮抗肌的抑制，
而无需任何信息传至脑。
脊髓与脑断开的动物，仍会在受到捏压时缩回爪子，
并仍会出现膝跳\omterm{def:g12:stretch-reflex:reflex}{反射}：脊髓是一个
中心，而不只是一根电缆。
\end{solution}

\begin{solution}{pb:g12:stretch-reflex:1}
\textbf{1.} $0.75 + 0.75 = \qty{1.5}{m}$。

\textbf{2.} 传导耗时 $30 - 4 - 4 - 1 = \qty{21}{ms}$：$1.5/0.021
\approx \qty{70}{m/s}$.

\textbf{3.} 往返各多约 \qty{10}{cm}，共 \qty{0.2}{m}：
$0.2/70 \approx \qty{3}{ms}$ 的额外时间，共约 \qty{33}{ms}。

\textbf{4.} 电信号即\omterm{def:g10:muscles-and-joints:muscle}{肌纤维}的激活；收缩本身，
即细丝的滑动以及对\omterm{def:g10:muscles-and-joints:muscle}{肌腱}的牵拉，
需要数十毫秒方能产生足够的力
以带动腿部运动。

\textbf{5.} 由脊髓传至脑的信息与由脑传回的指令，
以约 \qty{70}{m/s} 的速度传播 \qty{1.2}{m}，耗时 \qty{17}{ms}，再加上
若干\omterm{def:g12:stretch-reflex:synapse}{突触}及脑自身的处理时间，至少需 \qty{50}{ms}，
而做出决定所需的时间还要长得多。踢腿动作在任何脑信号
抵达运动\omterm{def:g12:stretch-reflex:neuron}{神经元}之前即已结束。

\textbf{6.} 以每秒 200 次计，\qty{10}{ms} 内有 2 个\omterm{prop:g12:stretch-reflex:actionpotential}{动作电位}
（静息时为 0.2 个）。频率发生变化；每个峰的大小与形状
则不变。

\textbf{7.} 相差 10 倍。

\textbf{8.} 以每秒 20 次计，递质剂量相隔 \qty{50}{ms}，
每一剂量在下一剂量到来前即已消退；运动\omterm{def:g12:stretch-reflex:neuron}{神经元}
始终未达到阈值。以每秒 200 次计，剂量相隔 \qty{5}{ms}，叠加
起来，越过阈值，运动\omterm{def:g12:stretch-reflex:neuron}{神经元}随之发放。

\textbf{9.} 达到阈值的运动\omterm{def:g12:stretch-reflex:neuron}{神经元}减少，且每个\omterm{def:g12:stretch-reflex:neuron}{神经元}发放的
\omterm{prop:g12:stretch-reflex:actionpotential}{动作电位}减少：收缩的运动单位减少，频率降低，
踢腿力度减弱。

\textbf{10.} \omterm{prop:g12:stretch-reflex:actionpotential}{动作电位}遵循“全或无”：其大小由
\omterm{def:g12:stretch-reflex:neuron}{神经元}决定，因而不携带关于刺激的信息。仅有
频率可以变化。

\textbf{11.} 乙酰胆碱；与受体结合后，它在
\omterm{def:g12:stretch-reflex:synapse}{突触}间隙中被一种酶在数毫秒内破坏。

\textbf{12.} \omterm{def:g10:muscles-and-joints:muscle}{肌纤维}上的受体：乙酰胆碱
已释放，但无法起作用。

\textbf{13.} 乙酰胆碱从囊泡中释放。在
实验室中，若将乙酰胆碱直接作用于肌肉，
在肉毒毒素（受体未被占据）作用下肌肉会收缩，而在箭毒
（受体被阻断）作用下则不会。

\textbf{14.} 感觉纤维的\omterm{prop:g12:stretch-reflex:actionpotential}{动作电位}；由于无信号传至脊髓，
运动\omterm{def:g12:stretch-reflex:neuron}{神经元}的\omterm{prop:g12:stretch-reflex:actionpotential}{动作电位}及肌电图均不出现；
但若在阻滞区以下刺激其神经，肌肉仍会收缩。

\textbf{15.} 箭毒阻断接收，肉毒毒素阻断释放，
麻醉剂阻断传导：三处阻断，同一块肌肉皆不收缩。

\textbf{16.} 感觉纤维还兴奋一个中间\omterm{def:g12:stretch-reflex:neuron}{神经元}，
该中间\omterm{def:g12:stretch-reflex:neuron}{神经元}抑制股后肌群的运动\omterm{def:g12:stretch-reflex:neuron}{神经元}，使其静息活动
受到抑制，拮抗肌则松弛至低于其静息张力。

\textbf{17.} 该\omterm{def:g12:stretch-reflex:reflex}{反射}只需脊髓，脊髓在该水平完好。
其\omterm{def:g12:stretch-reflex:reflex}{反射}亢进表明，脑通常发出
减弱\omterm{def:g12:stretch-reflex:reflex}{反射}的信号；缺少这些信号，\omterm{def:g12:stretch-reflex:reflex}{反射弧}便以全幅灵敏度运行。

\textbf{18.} 肌梭、感觉神经或背根、脊髓
在该节段的灰质、运动神经或腹根、肌肉
本身。直接刺激运动神经并记录肌肉，
可判断运动侧是否正常；在牵拉过程中记录感觉神经，
可判断感觉侧是否正常。

\textbf{19.} 一次叩击即可检验整个\omterm{def:g12:stretch-reflex:reflex}{反射弧}——感觉神经、
脊髓、运动神经、肌肉——以及脑对其的减弱作用：\omterm{def:g12:stretch-reflex:reflex}{反射}缺失、
减弱或亢进，即指向神经系统的某一层级，
且早于任何其他体征。

\textbf{20.} 约 \qty{70}{m/s}；肌梭\omterm{prop:g12:stretch-reflex:actionpotential}{动作电位}的
频率；乙酰胆碱与肌肉的
受体结合。
\end{solution}
